\subsection{Fast-carrier estimate at the illustrated amplitudes}\label{app:qif-estimate}
At $A_f=35$, the fast-carrier voltage excursion is
$A_f/(\tau_E\omega_c)\approx1.2$ per tone in model units, making these
runs finite-amplitude illustrations.
For two nearby carriers, the quadratic voltage term produces a constant
excitatory-drive shift $\delta\eta_E\simeq A_f^2/(\omega_c\tau_E)^2$:
about $1.5$ at $A_f=35$, compared with $0.12$ at $A_f=10$.
At $\bar\eta=0.85$, the Hopf threshold for a drive shift applied only to
the excitatory population is $\delta\eta_E^*\approx0.286$.
The larger field is therefore expected to cross the threshold; the smaller
one remains below it.

We refine this estimate using the full PING linear carrier responses and their
second-order DC source, solving for the carrier-averaged equilibrium and its
stability with the source evaluated at that equilibrium. At $\bar\eta=0.85$
and $A_f=35$, this gives $\delta\eta_E\approx1.65$ and a leading eigenvalue
with real part $+0.032$~ms$^{-1}$. This approximation isolates the mean shift.
Direct two-tone integrations retain the beat forcing and reproduce the off-beat
gamma line, providing a separate check.
